% Local Dirac notation commands for this chapter.
\providecommand{\ket}[1]{\left|#1\right\rangle}
\providecommand{\bra}[1]{\left\langle#1\right|}
\providecommand{\braket}[2]{\left\langle#1\middle|#2\right\rangle}
\providecommand{\mel}[3]{\left\langle#1\middle|#2\middle|#3\right\rangle}
\providecommand{\dd}{\mathrm{d}}
\providecommand{\ii}{\mathrm{i}}
\providecommand{\e}{\mathrm{e}}
\providecommand{\eps}{\varepsilon}
\providecommand{\unit}{\mathbbm{1}}

\chapter{Drehimpuls und Kugelflächenfunktionen}

\section{Ziel dieses Kapitels}

Der Drehimpuls ist in der Quantenmechanik eine der wichtigsten Observablen. Er tritt nicht nur bei der Bewegung eines Teilchens im Raum auf, sondern auch bei Spin-Systemen, Atomen, Auswahlregeln und Symmetrien. In diesem Kapitel geht es zunächst um den Bahndrehimpuls eines Teilchens in drei Dimensionen. Der Bahndrehimpuls ist der quantenmechanische Operator, der aus dem klassischen Ausdruck
\[
    \vec L = \vec r \times \vec p
\]
mittels kanonischer Quantisierung entsteht.

\begin{conceptbox}
Der zentrale Punkt ist: Drehimpuls ist der Generator von Rotationen. Wenn ein System rotationssymmetrisch ist, dann ist Drehimpuls erhalten. Für ein Zentralpotential hängt der Hamiltonoperator nur von \(r=|\vec r|\) ab. Dann kommutiert der Hamiltonoperator mit \(\vec L^{\,2}\) und mit \(L_z\). Dadurch können Energieeigenzustände nach Drehimpulsquantenzahlen klassifiziert werden.
\end{conceptbox}

Die wichtigsten Fragen dieses Kapitels sind:
\begin{itemize}
    \item Wie definiert man den Bahndrehimpulsoperator in der Quantenmechanik?
    \item Welche Kommutatorrelationen erfüllen \(L_x,L_y,L_z\)?
    \item Warum kann man nicht alle drei Drehimpulskomponenten gleichzeitig scharf messen?
    \item Warum kann man aber \(\vec L^{\,2}\) und eine Komponente, meistens \(L_z\), gleichzeitig diagonal wählen?
    \item Wie sehen die Operatoren in Kugelkoordinaten aus?
    \item Woher kommen die Quantenzahlen \(l\) und \(m\)?
    \item Was sind Kugelflächenfunktionen \(Y_l^m(\theta,\varphi)\)?
    \item Wie erhält man die Kugelflächenfunktionen für \(l=0\) und \(l=1\)?
    \item Warum sind beim Bahndrehimpuls nur ganzzahlige Werte von \(l\) möglich?
\end{itemize}

\section{Klassischer und quantenmechanischer Bahndrehimpuls}

\subsection{Klassischer Bahndrehimpuls}

In der klassischen Mechanik ist der Bahndrehimpuls eines Teilchens definiert durch
\[
    \vec L = \vec r \times \vec p.
\]
In Komponenten schreibt man
\[
    L_j = \sum_{k,\ell=1}^3 \eps_{jk\ell} x_k p_\ell,
\]
wobei \(\eps_{jk\ell}\) das Levi-Civita-Symbol ist. Explizit gilt
\[
    L_x = y p_z - z p_y,
    \qquad
    L_y = z p_x - x p_z,
    \qquad
    L_z = x p_y - y p_x.
\]

\begin{notebox}
Der Bahndrehimpuls misst nicht, wie schnell sich ein Teilchen um sich selbst dreht. Er beschreibt die Rotation der Bahnbewegung um einen Ursprung. Für ein Teilchen in einem Zentralpotential ist dieser Ursprung normalerweise das Kraftzentrum.
\end{notebox}

\subsection{Kanonische Quantisierung}

In der Ortsdarstellung ersetzt man
\[
    x_j \mapsto X_j=x_j,
    \qquad
    p_j \mapsto P_j=-\ii\hbar\frac{\partial}{\partial x_j}.
\]
Damit wird der klassische Bahndrehimpuls zum Operator
\begin{definitionbox}
\[
    L_j := \sum_{k,\ell=1}^3 \eps_{jk\ell} X_k P_\ell.
\]
In Ortsdarstellung gilt also
\[
    L_j = -\ii\hbar\sum_{k,\ell=1}^3 \eps_{jk\ell}x_k\frac{\partial}{\partial x_\ell}.
\]
\end{definitionbox}

Explizit erhält man
\begin{formulabox}
\[
\begin{aligned}
    L_x &= -\ii\hbar\left(y\frac{\partial}{\partial z}-z\frac{\partial}{\partial y}\right),\\
    L_y &= -\ii\hbar\left(z\frac{\partial}{\partial x}-x\frac{\partial}{\partial z}\right),\\
    L_z &= -\ii\hbar\left(x\frac{\partial}{\partial y}-y\frac{\partial}{\partial x}\right).
\end{aligned}
\]
\end{formulabox}

Das Quadrat des Gesamtdrehimpulses ist
\begin{definitionbox}
\[
    \vec L^{\,2}:=L_x^2+L_y^2+L_z^2.
\]
\end{definitionbox}

\section{Drehimpulskommutatoren}

\subsection{Kanonische Kommutatorrelationen als Ausgangspunkt}

Die Orts- und Impulsoperatoren erfüllen
\[
    [X_j,P_k]=\ii\hbar\delta_{jk},
    \qquad
    [X_j,X_k]=0,
    \qquad
    [P_j,P_k]=0.
\]
Diese Relationen reichen aus, um die Drehimpulsalgebra herzuleiten.

\begin{satzbox}
Die Komponenten des Bahndrehimpulses erfüllen
\[
    [L_i,L_j]=\ii\hbar\sum_{k=1}^3\eps_{ijk}L_k.
\]
Explizit:
\[
    [L_x,L_y]=\ii\hbar L_z,
    \qquad
    [L_y,L_z]=\ii\hbar L_x,
    \qquad
    [L_z,L_x]=\ii\hbar L_y.
\]
\end{satzbox}

\begin{proofbox}
Man verwendet
\[
    [AB,CD]=A[B,C]D+[A,C]DB+CA[B,D]+C[A,D]B.
\]
Mit
\[
    L_i=\eps_{iab}X_aP_b,
    \qquad
    L_j=\eps_{jcd}X_cP_d
\]
folgt
\[
\begin{aligned}
    [L_i,L_j]
    &=\eps_{iab}\eps_{jcd}[X_aP_b,X_cP_d]\\
    &=\eps_{iab}\eps_{jcd}\left(X_a[P_b,X_c]P_d+X_c[X_a,P_d]P_b\right)\\
    &=\eps_{iab}\eps_{jcd}\left(-\ii\hbar\delta_{bc}X_aP_d+\ii\hbar\delta_{ad}X_cP_b\right).
\end{aligned}
\]
Nach Umbenennung der Indizes und Verwendung der Identitäten des Levi-Civita-Symbols erhält man
\[
    [L_i,L_j]=\ii\hbar\eps_{ijk}L_k.
\]
\end{proofbox}

\subsection{Folgen für Messungen}

Da die Drehimpulskomponenten nicht miteinander kommutieren, können sie nicht alle gleichzeitig scharf gemessen werden. Insbesondere gilt
\[
    [L_x,L_y]=\ii\hbar L_z\neq 0.
\]
Daraus folgt eine Unschärferelation
\begin{formulabox}
\[
    \Delta L_x\,\Delta L_y\geq \frac{\hbar}{2}|\langle L_z\rangle|,
\]
und zyklisch analog.
\end{formulabox}

\begin{warningbox}
Man darf also nicht sagen: Ein Zustand hat gleichzeitig scharfe Werte von \(L_x\), \(L_y\) und \(L_z\). Standardmäßig wählt man \(\vec L^{\,2}\) und eine Komponente, meistens \(L_z\), als gleichzeitig messbare Observablen.
\end{warningbox}

\subsection{Kommutator von \(\vec L^{\,2}\) mit den Komponenten}

\begin{satzbox}
Für alle \(j\in\{x,y,z\}\) gilt
\[
    [\vec L^{\,2},L_j]=0.
\]
\end{satzbox}

\begin{proofbox}
Exemplarisch für \(L_z\):
\[
\begin{aligned}
    [\vec L^{\,2},L_z]
    &= [L_x^2+L_y^2+L_z^2,L_z]\\
    &= L_x[L_x,L_z]+[L_x,L_z]L_x
     + L_y[L_y,L_z]+[L_y,L_z]L_y.
\end{aligned}
\]
Mit
\[
    [L_x,L_z]=-\ii\hbar L_y,
    \qquad
    [L_y,L_z]=\ii\hbar L_x
\]
folgt
\[
\begin{aligned}
    [\vec L^{\,2},L_z]
    &= -\ii\hbar L_xL_y-\ii\hbar L_yL_x
       +\ii\hbar L_yL_x+\ii\hbar L_xL_y=0.
\end{aligned}
\]
Die anderen Komponenten folgen zyklisch.
\end{proofbox}

\begin{conceptbox}
Da \(\vec L^{\,2}\) mit jeder Komponente \(L_j\) kommutiert, kann man Zustände finden, die gleichzeitig Eigenzustände von \(\vec L^{\,2}\) und einer ausgewählten Komponente sind. Die Standardwahl ist \(\vec L^{\,2}\) und \(L_z\).
\end{conceptbox}

\section{Drehimpuls als Generator von Rotationen}

Eine Rotation im Hilbertraum wird durch einen unitären Operator dargestellt. Für eine infinitesimale Drehung um die Achse \(\vec n\) mit kleinem Winkel \(\alpha\) gilt
\[
    D(\alpha,\vec n)
    =\exp\left(-\frac{\ii}{\hbar}\alpha\,\vec n\cdot\vec L\right)
    \approx 1-\frac{\ii}{\hbar}\alpha\,\vec n\cdot\vec L.
\]

\begin{conceptbox}
Der Operator \(\vec n\cdot\vec L\) erzeugt Drehungen um die Achse \(\vec n\). Das ist analog zu
\[
    U(t)=\exp\left(-\frac{\ii}{\hbar}Ht\right),
\]
wobei der Hamiltonoperator die Zeitentwicklung erzeugt, und zu Translationen, die durch den Impuls erzeugt werden.
\end{conceptbox}

Wenn der Hamiltonoperator rotationsinvariant ist, dann gilt
\[
    [H,D(R)]=0
\]
für alle Rotationen \(R\). Für infinitesimale Rotationen folgt daraus
\[
    [H,L_j]=0.
\]
Dann ist der Drehimpuls erhalten:
\[
    \frac{\dd}{\dd t}\langle L_j\rangle
    =\frac{\ii}{\hbar}\langle[H,L_j]\rangle=0.
\]

\begin{examplebox}
Für ein Zentralpotential
\[
    H=\frac{\vec P^{\,2}}{2m}+V(r),
    \qquad r=\sqrt{x^2+y^2+z^2},
\]
ist der Hamiltonoperator rotationsinvariant. Daher kommutiert er mit \(\vec L^{\,2}\) und mit allen Komponenten von \(\vec L\). Deshalb kann man die Eigenzustände eines Zentralpotentials nach \(l\) und \(m\) klassifizieren.
\end{examplebox}

\section{Leiteroperatoren des Drehimpulses}

\subsection{Definition}

Analog zum harmonischen Oszillator definiert man Drehimpuls-Leiteroperatoren
\begin{definitionbox}
\[
    L_+ := L_x+\ii L_y,
    \qquad
    L_- := L_x-\ii L_y.
\]
\end{definitionbox}

Diese Operatoren verändern später die magnetische Quantenzahl \(m\), lassen aber \(l\) unverändert.

\subsection{Kommutatoren}

Aus der Drehimpulsalgebra folgt
\begin{formulabox}
\[
    [L_z,L_+]=\hbar L_+,
    \qquad
    [L_z,L_-]=-\hbar L_-.
\]
Außerdem gilt
\[
    [\vec L^{\,2},L_\pm]=0.
\]
\end{formulabox}

\begin{proofbox}
Zum Beispiel
\[
\begin{aligned}
    [L_z,L_+]
    &= [L_z,L_x+\ii L_y]\\
    &= [L_z,L_x]+\ii[L_z,L_y]\\
    &= \ii\hbar L_y+\ii(-\ii\hbar L_x)\\
    &= \hbar(L_x+\ii L_y)=\hbar L_+.
\end{aligned}
\]
Analog folgt \([L_z,L_-]=-\hbar L_-\).
\end{proofbox}

\subsection{Wirkung auf Eigenzustände}

Sei \(\ket{l,m}\) ein gemeinsamer Eigenzustand von \(\vec L^{\,2}\) und \(L_z\):
\[
    \vec L^{\,2}\ket{l,m}=\hbar^2 l(l+1)\ket{l,m},
    \qquad
    L_z\ket{l,m}=\hbar m\ket{l,m}.
\]
Dann gilt
\[
    L_z(L_+\ket{l,m})
    =([L_z,L_+]+L_+L_z)\ket{l,m}
    =\hbar(m+1)L_+\ket{l,m}.
\]
Also ist \(L_+\ket{l,m}\) ein Zustand mit magnetischer Quantenzahl \(m+1\). Ebenso senkt \(L_-\) den Wert von \(m\) um eins.

\begin{satzbox}
Die Leiteroperatoren wirken gemäß
\[
    L_+\ket{l,m}
    =\hbar\sqrt{l(l+1)-m(m+1)}\,\ket{l,m+1},
\]
\[
    L_-\ket{l,m}
    =\hbar\sqrt{l(l+1)-m(m-1)}\,\ket{l,m-1}.
\]
\end{satzbox}

\begin{derivationbox}
Man benutzt
\[
    L_-L_+ = \vec L^{\,2}-L_z^2-\hbar L_z,
    \qquad
    L_+L_- = \vec L^{\,2}-L_z^2+\hbar L_z.
\]
Dann ist
\[
\begin{aligned}
    \|L_+\ket{l,m}\|^2
    &=\mel{l,m}{L_-L_+}{l,m}\\
    &=\hbar^2\left[l(l+1)-m^2-m\right]\\
    &=\hbar^2\left[l(l+1)-m(m+1)\right].
\end{aligned}
\]
Die Wurzel dieses Ausdrucks ist der Normierungsfaktor. Analog erhält man die Formel für \(L_-\).
\end{derivationbox}

\begin{examtipbox}
Die Formeln für \(L_\pm\ket{l,m}\) sind Standardformeln. In Aufgaben muss man sie entweder direkt anwenden oder aus den Normen \(\|L_\pm\ket{l,m}\|^2\) herleiten können.
\end{examtipbox}

\section{Quantenzahlen \(l\) und \(m\)}

\subsection{Allgemeine Drehimpulsquantisierung}

Aus der Algebra folgt allgemein:
\[
    l=0,\frac12,1,\frac32,2,\dots,
    \qquad
    m=-l,-l+1,\dots,l-1,l.
\]
Für einen gegebenen Wert von \(l\) gibt es also
\[
    2l+1
\]
verschiedene Werte von \(m\). Diese Zustände spannen eine irreduzible Darstellung der Drehimpulsalgebra auf.

\begin{notebox}
Diese allgemeine Aussage gilt für jeden quantenmechanischen Drehimpuls, also auch für Spin. Für den Bahndrehimpuls in Ortsdarstellung kommen jedoch nur ganzzahlige Werte \(l=0,1,2,\dots\) vor. Halbzahlige Werte gehören zum Spin, nicht zu gewöhnlichen skalaren Wellenfunktionen auf dem Raum.
\end{notebox}

\subsection{Warum ist \(m\) beschränkt?}

Da Normen nicht negativ sein können, gilt
\[
    \|L_+\ket{l,m}\|^2\geq 0,
    \qquad
    \|L_-\ket{l,m}\|^2\geq 0.
\]
Also
\[
    l(l+1)-m(m+1)\geq 0,
    \qquad
    l(l+1)-m(m-1)\geq 0.
\]
Daraus folgt
\[
    -l\leq m\leq l.
\]
Am oberen Ende muss gelten
\[
    L_+\ket{l,l}=0,
\]
am unteren Ende
\[
    L_-\ket{l,-l}=0.
\]

\begin{conceptbox}
Die Quantenzahl \(l\) bestimmt die Länge des Drehimpulsvektors über \(\hbar^2l(l+1)\). Die Quantenzahl \(m\) bestimmt die Projektion auf die \(z\)-Achse über \(\hbar m\). Der Betrag der Projektion kann nie größer sein als die gesamte Drehimpulslänge.
\end{conceptbox}

\section{Drehimpulsoperatoren in Kugelkoordinaten}

\subsection{Kugelkoordinaten}

Wir verwenden
\[
    x=r\sin\theta\cos\varphi,
    \qquad
    y=r\sin\theta\sin\varphi,
    \qquad
    z=r\cos\theta.
\]
Dabei ist
\[
    r\geq 0,
    \qquad
    0\leq \theta\leq \pi,
    \qquad
    0\leq \varphi<2\pi.
\]
Der Winkel \(\theta\) ist der Polarwinkel von der positiven \(z\)-Achse aus. Der Winkel \(\varphi\) ist der Azimutwinkel in der \(xy\)-Ebene.

\begin{warningbox}
In manchen Skripten wird \(\varphi\) und \(\theta\) anders benannt. In diesem Kapitel gilt immer: \(\theta\) ist der Polarwinkel und \(\varphi\) der Azimutwinkel.
\end{warningbox}

\subsection{Der Operator \(L_z\)}

In Kugelkoordinaten vereinfacht sich \(L_z\) besonders stark:
\begin{formulabox}
\[
    L_z=-\ii\hbar\frac{\partial}{\partial\varphi}.
\]
\end{formulabox}

\begin{derivationbox}
Aus
\[
    L_z=-\ii\hbar\left(x\frac{\partial}{\partial y}-y\frac{\partial}{\partial x}\right)
\]
und der Kettenregel folgt, dass dieser Operator genau Änderungen des Azimutwinkels erzeugt. Eine Rotation um die \(z\)-Achse verändert nur \(\varphi\), nicht \(r\) und nicht \(\theta\). Daher ist
\[
    L_z=-\ii\hbar\partial_\varphi.
\]
\end{derivationbox}

\subsection{Die Leiteroperatoren in Kugelkoordinaten}

Man erhält
\begin{formulabox}
\[
    L_\pm
    =\hbar e^{\pm\ii\varphi}
    \left(\pm\frac{\partial}{\partial\theta}
    +\ii\cot\theta\frac{\partial}{\partial\varphi}\right).
\]
\end{formulabox}

\subsection{Das Quadrat des Drehimpulses}

Für das Drehimpulsquadrat gilt in Kugelkoordinaten
\begin{formulabox}
\[
    \vec L^{\,2}
    =-\hbar^2\left[
    \frac{1}{\sin\theta}\frac{\partial}{\partial\theta}
    \left(\sin\theta\frac{\partial}{\partial\theta}\right)
    +\frac{1}{\sin^2\theta}\frac{\partial^2}{\partial\varphi^2}
    \right].
\]
\end{formulabox}

\begin{notebox}
Auffällig ist: In \(\vec L^{\,2}\), \(L_z\) und \(L_\pm\) kommt keine Ableitung nach \(r\) vor. Der Bahndrehimpuls wirkt nur auf die Winkelabhängigkeit einer Wellenfunktion. Der radiale Anteil bleibt davon unberührt.
\end{notebox}

\section{Eigenwertproblem von \(\vec L^{\,2}\) und \(L_z\)}

\subsection{Gleichzeitige Eigenfunktionen}

Da
\[
    [\vec L^{\,2},L_z]=0,
\]
kann man gemeinsame Eigenfunktionen wählen:
\[
    \vec L^{\,2}\psi(r,\theta,\varphi)
    =\hbar^2l(l+1)\psi(r,\theta,\varphi),
\]
\[
    L_z\psi(r,\theta,\varphi)
    =\hbar m\psi(r,\theta,\varphi).
\]
Da die Operatoren nur auf die Winkel wirken, verwendet man den Separationsansatz
\[
    \psi(r,\theta,\varphi)=R(r)Y_l^m(\theta,\varphi).
\]

\begin{definitionbox}
Die Funktionen \(Y_l^m(\theta,\varphi)\) heißen Kugelflächenfunktionen. Sie sind die gemeinsamen Eigenfunktionen von \(\vec L^{\,2}\) und \(L_z\) auf der Einheitssphäre:
\[
    \vec L^{\,2}Y_l^m=\hbar^2l(l+1)Y_l^m,
    \qquad
    L_zY_l^m=\hbar mY_l^m.
\]
\end{definitionbox}

\begin{warningbox}
Die Operatoren \(\vec L^{\,2}\) und \(L_z\) bestimmen nur den Winkelanteil. Der Radialanteil \(R(r)\) wird durch den konkreten Hamiltonoperator bestimmt, zum Beispiel durch ein Zentralpotential. Deshalb bilden \(\vec L^{\,2}\) und \(L_z\) allein keinen vollständigen Satz kommutierender Observablen für den gesamten Hilbertraum \(L^2(\mathbb R^3)\).
\end{warningbox}

\subsection{Differentialgleichung für \(L_z\)}

Aus
\[
    L_zY_l^m=\hbar mY_l^m
\]
folgt
\[
    -\ii\hbar\frac{\partial}{\partial\varphi}Y_l^m(\theta,\varphi)
    =\hbar mY_l^m(\theta,\varphi).
\]
Nach Division durch \(\hbar\):
\[
    -\ii\frac{\partial}{\partial\varphi}Y_l^m=mY_l^m.
\]
Die Lösung hat die Form
\[
    Y_l^m(\theta,\varphi)=F_l^m(\theta)e^{\ii m\varphi}.
\]
Da eine skalare Wellenfunktion bei \(\varphi\mapsto \varphi+2\pi\) eindeutig sein muss, fordert man
\[
    Y_l^m(\theta,\varphi+2\pi)=Y_l^m(\theta,\varphi).
\]
Damit gilt
\[
    e^{\ii m2\pi}=1.
\]
Also muss \(m\) ganzzahlig sein:
\begin{formulabox}
\[
    m\in\mathbb Z.
\]
\end{formulabox}

\begin{conceptbox}
Für den Bahndrehimpuls sind \(m\) und damit auch \(l\) ganzzahlig. Halbzahlige Drehimpulse treten bei Spinoren auf. Eine Spin-\(1/2\)-Wellenfunktion darf bei einer Drehung um \(2\pi\) ein Vorzeichen wechseln; eine gewöhnliche skalare Ortswellenfunktion darf das nicht.
\end{conceptbox}

\subsection{Differentialgleichung für \(\vec L^{\,2}\)}

Setzt man die explizite Form von \(\vec L^{\,2}\) ein, erhält man
\begin{formulabox}
\[
    -\left[
    \frac{1}{\sin\theta}\frac{\partial}{\partial\theta}
    \left(\sin\theta\frac{\partial}{\partial\theta}\right)
    +\frac{1}{\sin^2\theta}\frac{\partial^2}{\partial\varphi^2}
    \right]Y_l^m
    =l(l+1)Y_l^m.
\]
\end{formulabox}

Mit
\[
    Y_l^m(\theta,\varphi)=F_l^m(\theta)e^{\ii m\varphi}
\]
folgt
\[
    -\frac{1}{\sin\theta}\frac{\dd}{\dd\theta}
    \left(\sin\theta\frac{\dd F_l^m}{\dd\theta}\right)
    +\frac{m^2}{\sin^2\theta}F_l^m
    =l(l+1)F_l^m.
\]
Das ist die Winkelgleichung für die Kugelflächenfunktionen.

\section{Explizite Kugelflächenfunktionen}

\subsection{Allgemeine Form}

Die normierten Kugelflächenfunktionen können geschrieben werden als
\begin{formulabox}
\[
    Y_l^m(\theta,\varphi)
    =(-1)^m\sqrt{\frac{2l+1}{4\pi}\frac{(l-m)!}{(l+m)!}}
    P_l^m(\cos\theta)e^{\ii m\varphi}
\]
für \(m\geq 0\). Für negative \(m\) verwendet man
\[
    Y_l^{-m}(\theta,\varphi)=(-1)^m\left(Y_l^m(\theta,\varphi)\right)^*.
\]
Hier sind \(P_l^m\) die zugeordneten Legendre-Polynome.
\end{formulabox}

\begin{notebox}
In der Literatur gibt es verschiedene Phasenkonventionen. Die obige Konvention enthält den sogenannten Condon-Shortley-Phasenfaktor \((-1)^m\). Physikalische Wahrscheinlichkeiten ändern sich durch eine globale Phasenkonvention nicht.
\end{notebox}

\subsection{Normierung und Orthogonalität}

Auf der Einheitssphäre gilt das Flächenelement
\[
    \dd\Omega=\sin\theta\,\dd\theta\,\dd\varphi.
\]
Die Kugelflächenfunktionen erfüllen
\begin{formulabox}
\[
    \int_0^{2\pi}\dd\varphi\int_0^\pi\dd\theta\,\sin\theta\,
    \left(Y_l^m(\theta,\varphi)\right)^*Y_{l'}^{m'}(\theta,\varphi)
    =\delta_{ll'}\delta_{mm'}.
\]
\end{formulabox}

Außerdem bilden sie eine vollständige Basis für quadratintegrierbare Funktionen auf der Kugeloberfläche:
\[
    f(\theta,\varphi)=\sum_{l=0}^{\infty}\sum_{m=-l}^{l}c_{lm}Y_l^m(\theta,\varphi).
\]
Die Koeffizienten sind
\[
    c_{lm}=\int \dd\Omega\,\left(Y_l^m\right)^*f.
\]

\section{Die wichtigsten Kugelflächenfunktionen}

\subsection{Der Fall \(l=0\)}

Für \(l=0\) ist nur \(m=0\) erlaubt. Die Kugelflächenfunktion ist konstant:
\begin{formulabox}
\[
    Y_0^0(\theta,\varphi)=\frac{1}{\sqrt{4\pi}}.
\]
\end{formulabox}

\begin{conceptbox}
Der Zustand \(l=0\) ist kugelsymmetrisch. Er besitzt keine bevorzugte Richtung. In Atomphysik spricht man von einem \(s\)-Zustand.
\end{conceptbox}

\subsection{Der Fall \(l=1\)}

Für \(l=1\) gibt es drei Werte:
\[
    m=-1,0,1.
\]
Eine übliche normierte Wahl ist
\begin{formulabox}
\[
    Y_1^1(\theta,\varphi)
    =-\sqrt{\frac{3}{8\pi}}\sin\theta\,e^{\ii\varphi},
\]
\[
    Y_1^0(\theta,\varphi)
    =\sqrt{\frac{3}{4\pi}}\cos\theta,
\]
\[
    Y_1^{-1}(\theta,\varphi)
    =\sqrt{\frac{3}{8\pi}}\sin\theta\,e^{-\ii\varphi}.
\]
\end{formulabox}

\begin{notebox}
Die drei \(l=1\)-Zustände nennt man oft \(p\)-Zustände. Sie entsprechen drei verschiedenen magnetischen Quantenzahlen. Reelle Linearkombinationen dieser komplexen Kugelflächenfunktionen ergeben die bekannten \(p_x\)-, \(p_y\)- und \(p_z\)-artigen Winkelabhängigkeiten.
\end{notebox}

\section{Herleitung der \(l=1\)-Kugelflächenfunktionen mit Leiteroperatoren}

\subsection{Start beim höchsten Gewicht}

Für den höchsten Zustand gilt
\[
    L_+Y_l^l=0.
\]
Für \(l=1\) setzen wir
\[
    Y_1^1(\theta,\varphi)=F(\theta)e^{\ii\varphi}.
\]
Dann gilt
\[
    L_+Y_1^1
    =\hbar e^{\ii\varphi}\left(\frac{\partial}{\partial\theta}
    +\ii\cot\theta\frac{\partial}{\partial\varphi}\right)
    \left[F(\theta)e^{\ii\varphi}\right].
\]
Da
\[
    \frac{\partial}{\partial\varphi}e^{\ii\varphi}=\ii e^{\ii\varphi},
\]
erhält man
\[
    L_+Y_1^1
    =\hbar e^{2\ii\varphi}\left(F'(\theta)-\cot\theta F(\theta)\right).
\]
Die Bedingung \(L_+Y_1^1=0\) liefert
\[
    F'(\theta)-\cot\theta F(\theta)=0.
\]
Also
\[
    \frac{F'}{F}=\cot\theta.
\]
Integration ergibt
\[
    \ln F=\ln(\sin\theta)+\text{const.}
\]
und damit
\[
    F(\theta)=C\sin\theta.
\]
Somit
\[
    Y_1^1(\theta,\varphi)=C\sin\theta e^{\ii\varphi}.
\]
Mit der Standardphasenkonvention und Normierung ist
\[
    C=-\sqrt{\frac{3}{8\pi}}.
\]

\subsection{Absenken mit \(L_-\)}

Die Leiterformel liefert für \(l=1,m=1\):
\[
    L_-Y_1^1
    =\hbar\sqrt{1(1+1)-1(1-1)}Y_1^0
    =\hbar\sqrt2\,Y_1^0.
\]
Also
\[
    Y_1^0=\frac{1}{\hbar\sqrt2}L_-Y_1^1.
\]
Man kann damit aus \(Y_1^1\) die Funktion \(Y_1^0\) berechnen. Ein weiteres Anwenden von \(L_-\) liefert \(Y_1^{-1}\):
\[
    Y_1^{-1}=\frac{1}{\hbar\sqrt2}L_-Y_1^0.
\]
So erhält man genau die drei oben angegebenen Funktionen.

\begin{examtipbox}
Für Übungsblatt- oder Tafelaufgaben ist diese Methode sehr wichtig: Man startet mit \(L_+Y_l^l=0\), löst die einfache Differentialgleichung für den höchsten Zustand und erzeugt die restlichen Zustände durch \(L_-\).
\end{examtipbox}

\section{Überprüfung der Eigenwertgleichungen für \(l=1\)}

\subsection{Überprüfung von \(L_z\)}

Für \(Y_1^1\) gilt
\[
    Y_1^1=-\sqrt{\frac{3}{8\pi}}\sin\theta e^{\ii\varphi}.
\]
Dann ist
\[
    L_zY_1^1
    =-\ii\hbar\frac{\partial}{\partial\varphi}Y_1^1
    =-\ii\hbar(\ii)Y_1^1
    =\hbar Y_1^1.
\]
Also ist \(m=1\). Für \(Y_1^0\) gibt es keine \(\varphi\)-Abhängigkeit, daher
\[
    L_zY_1^0=0.
\]
Für \(Y_1^{-1}\propto e^{-\ii\varphi}\) gilt
\[
    L_zY_1^{-1}=-\hbar Y_1^{-1}.
\]

\subsection{Überprüfung von \(\vec L^{\,2}\)}

Für alle drei \(l=1\)-Funktionen muss gelten
\[
    \vec L^{\,2}Y_1^m=\hbar^2\cdot 1\cdot 2\,Y_1^m=2\hbar^2Y_1^m.
\]
Für \(m=0\) ist die Rechnung besonders kurz. Es gilt
\[
    Y_1^0=C\cos\theta.
\]
Da keine \(\varphi\)-Abhängigkeit vorhanden ist,
\[
    \vec L^{\,2}Y_1^0
    =-\hbar^2\frac{1}{\sin\theta}\frac{\partial}{\partial\theta}
    \left(\sin\theta\frac{\partial}{\partial\theta}C\cos\theta\right).
\]
Nun
\[
    \frac{\partial}{\partial\theta}C\cos\theta=-C\sin\theta,
\]
also
\[
    \sin\theta\frac{\partial}{\partial\theta}C\cos\theta=-C\sin^2\theta.
\]
Weiter
\[
    \frac{\partial}{\partial\theta}(-C\sin^2\theta)=-2C\sin\theta\cos\theta.
\]
Damit
\[
    \vec L^{\,2}Y_1^0
    =-\hbar^2\frac{1}{\sin\theta}(-2C\sin\theta\cos\theta)
    =2\hbar^2C\cos\theta
    =2\hbar^2Y_1^0.
\]

\begin{answerbox}
Die Funktionen \(Y_1^1,Y_1^0,Y_1^{-1}\) sind gemeinsame Eigenfunktionen von \(\vec L^{\,2}\) und \(L_z\) mit
\[
    l=1,
    \qquad
    m=1,0,-1.
\]
\end{answerbox}

\section{Winkelanteil und Radialanteil in Zentralpotentialen}

\subsection{Hamiltonoperator bei Zentralpotentialen}

Für ein Zentralpotential
\[
    V(\vec r)=V(r)
\]
lautet der Hamiltonoperator
\[
    H=-\frac{\hbar^2}{2m}\nabla^2+V(r).
\]
In Kugelkoordinaten kann der Laplace-Operator geschrieben werden als
\[
    \nabla^2
    =\frac{1}{r^2}\frac{\partial}{\partial r}
    \left(r^2\frac{\partial}{\partial r}\right)
    -\frac{\vec L^{\,2}}{\hbar^2r^2}.
\]
Damit wird
\begin{formulabox}
\[
    H=-\frac{\hbar^2}{2m}\frac{1}{r^2}\frac{\partial}{\partial r}
    \left(r^2\frac{\partial}{\partial r}\right)
    +\frac{\vec L^{\,2}}{2mr^2}+V(r).
\]
\end{formulabox}

\begin{conceptbox}
Der Term
\[
    \frac{\vec L^{\,2}}{2mr^2}
\]
wirkt wie eine zusätzliche effektive Energie, die aus der Winkelbewegung stammt. In der radialen Schrödingergleichung führt er zum Zentrifugalterm.
\end{conceptbox}

\subsection{Separation der Wellenfunktion}

Da \(H\), \(\vec L^{\,2}\) und \(L_z\) miteinander kommutieren, kann man Eigenfunktionen der Form
\[
    \psi_{nlm}(r,\theta,\varphi)=R_{nl}(r)Y_l^m(\theta,\varphi)
\]
wählen. Dabei bestimmt \(Y_l^m\) den Winkelanteil und \(R_{nl}(r)\) den Radialanteil.

\begin{notebox}
Das ist die Grundlage der quantenmechanischen Beschreibung des Wasserstoffatoms und vieler dreidimensionaler Potentialprobleme. Die Kugelflächenfunktionen sind dabei universell; nur die Radialfunktionen hängen vom konkreten Potential ab.
\end{notebox}

\section{Parität der Kugelflächenfunktionen}

Die Raumspiegelung wirkt auf Ortsvektoren als
\[
    \vec r\mapsto -\vec r.
\]
In Kugelkoordinaten entspricht das
\[
    (\theta,\varphi)\mapsto (\pi-\theta,\varphi+\pi).
\]
Für Kugelflächenfunktionen gilt
\begin{formulabox}
\[
    Y_l^m(\pi-\theta,\varphi+\pi)=(-1)^lY_l^m(\theta,\varphi).
\]
\end{formulabox}

\begin{conceptbox}
Die Parität eines Bahndrehimpulszustands ist \((-1)^l\). Zustände mit geradem \(l\) haben gerade Parität, Zustände mit ungeradem \(l\) haben ungerade Parität.
\end{conceptbox}

\section{Zusammenhang mit Drehimpulsdarstellungen}

Die Zustände \(\ket{l,m}\) bilden für festes \(l\) einen \((2l+1)\)-dimensionalen Vektorraum. In diesem Raum sind \(\vec L^{\,2}\) und \(L_z\) diagonal:
\[
    \vec L^{\,2}\ket{l,m}=\hbar^2l(l+1)\ket{l,m},
    \qquad
    L_z\ket{l,m}=\hbar m\ket{l,m}.
\]
Die Operatoren \(L_x\) und \(L_y\) sind nicht diagonal in dieser Basis, können aber aus \(L_\pm\) rekonstruiert werden:
\[
    L_x=\frac12(L_++L_-),
    \qquad
    L_y=\frac{1}{2\ii}(L_+-L_-).
\]

\subsection{Beispiel: Matrixdarstellung für \(l=1\)}

In der Basis
\[
    \{\ket{1,1},\ket{1,0},\ket{1,-1}\}
\]
ist
\begin{formulabox}
\[
    L_z=\hbar
    \begin{pmatrix}
        1&0&0\\
        0&0&0\\
        0&0&-1
    \end{pmatrix}.
\]
\end{formulabox}

Die Leiteroperatoren sind
\[
    L_+\ket{1,-1}=\hbar\sqrt2\ket{1,0},
    \qquad
    L_+\ket{1,0}=\hbar\sqrt2\ket{1,1},
    \qquad
    L_+\ket{1,1}=0.
\]
Daher
\[
    L_+=\hbar
    \begin{pmatrix}
        0&\sqrt2&0\\
        0&0&\sqrt2\\
        0&0&0
    \end{pmatrix},
    \qquad
    L_-=\hbar
    \begin{pmatrix}
        0&0&0\\
        \sqrt2&0&0\\
        0&\sqrt2&0
    \end{pmatrix}.
\]
Damit folgt
\begin{formulabox}
\[
    L_x=\frac{\hbar}{\sqrt2}
    \begin{pmatrix}
        0&1&0\\
        1&0&1\\
        0&1&0
    \end{pmatrix},
    \qquad
    L_y=\frac{\hbar}{\sqrt2}
    \begin{pmatrix}
        0&-\ii&0\\
        \ii&0&-\ii\\
        0&\ii&0
    \end{pmatrix}.
\]
\end{formulabox}

\section{Typische Aufgaben und Rechenmuster}

\subsection{Aufgabe: Ist ein Zustand Eigenzustand von \(L_z\)?}

\begin{examplebox}
Gegeben sei
\[
    f(\theta,\varphi)=C\sin\theta e^{2\ii\varphi}.
\]
Dann
\[
    L_z f=-\ii\hbar\frac{\partial}{\partial\varphi}f
    =-\ii\hbar(2\ii)f=2\hbar f.
\]
Also ist \(f\) ein Eigenzustand von \(L_z\) mit \(m=2\). Ob \(f\) auch Eigenzustand von \(\vec L^{\,2}\) ist, muss separat geprüft werden.
\end{examplebox}

\subsection{Aufgabe: Leiteroperator anwenden}

\begin{examplebox}
Gesucht ist \(L_-\ket{2,1}\). Man benutzt
\[
    L_-\ket{l,m}
    =\hbar\sqrt{l(l+1)-m(m-1)}\ket{l,m-1}.
\]
Einsetzen von \(l=2,m=1\) liefert
\[
    L_-\ket{2,1}
    =\hbar\sqrt{2\cdot 3-1\cdot 0}\ket{2,0}
    =\hbar\sqrt6\ket{2,0}.
\]
\end{examplebox}

\subsection{Aufgabe: Erwartungswert von \(L_z\)}

Sei
\[
    \ket{\psi}=a\ket{l,m_1}+b\ket{l,m_2},
    \qquad
    |a|^2+|b|^2=1.
\]
Dann ist
\[
    \langle L_z\rangle
    =|a|^2\hbar m_1+|b|^2\hbar m_2,
\]
weil \(L_z\) in der \(\ket{l,m}\)-Basis diagonal ist.

\subsection{Aufgabe: Wahrscheinlichkeiten bei Messung von \(L_z\)}

Ist ein Zustand als Superposition
\[
    \ket{\psi}=\sum_m c_m\ket{l,m}
\]
gegeben, dann ist die Wahrscheinlichkeit für das Messergebnis \(\hbar m\)
\[
    P(m)=|c_m|^2.
\]
Nach der Messung befindet sich das System im zugehörigen Eigenzustand \(\ket{l,m}\), sofern keine Entartung berücksichtigt werden muss.

\section{Häufige Fehler}

\begin{warningbox}
\textbf{Fehler 1:} \(l\) und \(m\) verwechseln. Der Eigenwert von \(\vec L^{\,2}\) ist \(\hbar^2l(l+1)\), nicht \(\hbar^2m(m+1)\). Der Eigenwert von \(L_z\) ist \(\hbar m\).
\end{warningbox}

\begin{warningbox}
\textbf{Fehler 2:} Zu glauben, dass \(L_x,L_y,L_z\) gleichzeitig messbar sind. Das stimmt nicht, weil die Komponenten nicht kommutieren.
\end{warningbox}

\begin{warningbox}
\textbf{Fehler 3:} Den Radialanteil durch Drehimpulsoperatoren bestimmen zu wollen. Die Drehimpulsoperatoren wirken nur auf Winkel. Der Radialanteil kommt aus dem Hamiltonoperator.
\end{warningbox}

\begin{warningbox}
\textbf{Fehler 4:} Beim Bahndrehimpuls halbzahlige \(l\)-Werte zulassen. Für gewöhnliche skalare Ortswellenfunktionen sind nur \(l=0,1,2,\dots\) möglich. Halbzahlige Werte gehören zu Spin-Darstellungen.
\end{warningbox}

\section{Prüfungsrelevante Zusammenfassung}

\begin{examtipbox}
Für Prüfung und Proseminar sollte man sicher können:
\begin{itemize}
    \item Definition des Bahndrehimpulses:
    \[
        L_j=\sum_{k,\ell}\eps_{jk\ell}X_kP_\ell.
    \]
    \item Kommutatorrelationen:
    \[
        [L_i,L_j]=\ii\hbar\eps_{ijk}L_k,
        \qquad
        [\vec L^{\,2},L_j]=0.
    \]
    \item Bedeutung simultaner Messbarkeit: \(\vec L^{\,2}\) und \(L_z\) kommutieren, aber \(L_x,L_y,L_z\) nicht paarweise.
    \item Leiteroperatoren:
    \[
        L_\pm=L_x\pm\ii L_y.
    \]
    \item Wirkung:
    \[
        L_\pm\ket{l,m}
        =\hbar\sqrt{l(l+1)-m(m\pm1)}\ket{l,m\pm1}.
    \]
    \item Kugelkoordinatenformeln:
    \[
        L_z=-\ii\hbar\partial_\varphi,
    \]
    \[
        \vec L^{\,2}
        =-\hbar^2\left[
        \frac{1}{\sin\theta}\partial_\theta(\sin\theta\partial_\theta)
        +\frac{1}{\sin^2\theta}\partial_\varphi^2
        \right].
    \]
    \item Eigenwertgleichungen:
    \[
        \vec L^{\,2}Y_l^m=\hbar^2l(l+1)Y_l^m,
        \qquad
        L_zY_l^m=\hbar mY_l^m.
    \]
    \item Erlaubte Werte beim Bahndrehimpuls:
    \[
        l=0,1,2,\dots,
        \qquad
        m=-l,-l+1,\dots,l.
    \]
    \item Die expliziten Kugelflächenfunktionen für \(l=0\) und \(l=1\).
\end{itemize}
\end{examtipbox}

\section{Kurzantworten}

\begin{answerbox}
\textbf{Was ist Bahndrehimpuls?}
\[
    \vec L=\vec r\times\vec p,
    \qquad
    L_j=\eps_{jk\ell}X_kP_\ell.
\]
Er beschreibt die Drehbewegung der Ortswellenfunktion um einen Ursprung.
\end{answerbox}

\begin{answerbox}
\textbf{Warum wählt man \(\vec L^{\,2}\) und \(L_z\)?}
Weil
\[
    [\vec L^{\,2},L_z]=0,
\]
aber
\[
    [L_x,L_y]=\ii\hbar L_z\neq 0.
\]
Man kann also \(\vec L^{\,2}\) und eine Komponente gleichzeitig scharf messen, aber nicht alle drei Komponenten.
\end{answerbox}

\begin{answerbox}
\textbf{Was sind Kugelflächenfunktionen?}
Kugelflächenfunktionen \(Y_l^m(\theta,\varphi)\) sind die gemeinsamen Eigenfunktionen von \(\vec L^{\,2}\) und \(L_z\):
\[
    \vec L^{\,2}Y_l^m=\hbar^2l(l+1)Y_l^m,
    \qquad
    L_zY_l^m=\hbar mY_l^m.
\]
Sie bilden die natürliche Winkelbasis für dreidimensionale rotationssymmetrische Probleme.
\end{answerbox}

\begin{answerbox}
\textbf{Was sind die erlaubten Quantenzahlen?}
Für den Bahndrehimpuls gilt
\[
    l=0,1,2,\dots,
    \qquad
    m=-l,-l+1,\dots,l.
\]
Für jedes \(l\) gibt es \(2l+1\) verschiedene \(m\)-Zustände.
\end{answerbox}
